% GENERATED FILE. Edit canonical lessons, then run npm run generate:books.
% canonical-source-hash: fnv1a64:91ff9459bfd33b5e
% canonical-lessons: JA-C222-nuru, JA-C222-omimai, JA-C222-kaifuku, JA-C222-daijoubu, JA-C222-gaman

\chapter{Medicine, Visits and Getting Well}
\label{ch:ja-medicine-visits-and-getting-well}
\hlchaptermodality{222}

\begin{chapteropening}
\textbf{By the end of this chapter:} \emph{I can talk about medicine, sick visits and recovery.}

Complete the last lesson of chapter 222: I can talk about medicine, sick visits and recovery.
\end{chapteropening}

\section[\texorpdfstring{\ja{ぬる} (nuru)}{ぬる (nuru)}]{\ja{ぬる} (nuru) — to apply (cream); to spread on}

\label{lesson:JA-C222-nuru}



\begin{quote}
Before the new one: say the Japanese for treatment, then the Japanese for prevention.
\end{quote}

\subsection*{You'll want to know: \ja{ぬる}}
\textbf{\ja{ぬる}} — \emph{nuru} — \textquotedblleft{}to apply (cream); to spread on\textquotedblright{}.

\begin{quote}
\textbf{\ja{きずに} \ja{くすりを} \ja{ぬります。}} — \emph{kizu ni kusuri o nurimasu} — \textquotedblleft{}I put ointment on the wound.\textquotedblright{}
\end{quote}

Polite \textbf{\ja{ぬります}}, \emph{nurimasu}; a godan verb.

\begin{cousinweb}
\textbf{\ja{ぬる}} (\emph{nuru}) is a native Japanese word (wago), inherited rather than borrowed from Chinese; its kanji was attached later as a kun'yomi, and its older origin is not settled (JMdict, KANJIDIC2).
\end{cousinweb}

\begin{grammarlens}[title={What you've built}]
One more word for everyday life at A2.
\end{grammarlens}

\subsection*{Your turn}
\begin{itemize}
\raggedright
  \item \emph{Say it:} \emph{nuru}
  \item \emph{Say it:} \emph{nuru}, once more
  \item \emph{Say it:} \emph{kizu ni kusuri o nurimasu} — I put ointment on the wound.
  \item \emph{Recall:} say the Japanese for treatment, then the Japanese for prevention, then say \emph{nuru} again
\end{itemize}

\begin{tcolorbox}[breakable,colback=teal!4,colframe=teal!35!black,title={Before you move on}]
What does \ja{ぬる} mean? (\textquotedblleft{}To apply (cream); to spread on\textquotedblright{}.) Say it once more.
\end{tcolorbox}
\section[\texorpdfstring{\ja{おみまい} (omimai)}{おみまい (omimai)}]{\ja{おみまい} (omimai) — a sick visit; calling on someone who is ill}

\label{lesson:JA-C222-omimai}



\begin{quote}
Before the new one: say the Japanese for prevention, then the Japanese for to apply.
\end{quote}

\subsection*{You'll want to know: \ja{おみまい}}
\textbf{\ja{おみまい}} — \emph{omimai} — \textquotedblleft{}a sick visit; calling on someone who is ill\textquotedblright{}.

\begin{quote}
\textbf{\ja{ともだちの} \ja{おみまいに} \ja{いきませんか。}} — \emph{tomodachi no omimai ni ikimasen ka} — \textquotedblleft{}Shall we go and visit our sick friend?\textquotedblright{}
\end{quote}

\begin{cousinweb}
\textbf{\ja{おみまい}} (\emph{omimai}) is native Japanese (wago), built from polite \emph{o-} and \emph{mimai}, \textquotedblleft{}calling on someone\textquotedblright{}, inherited rather than borrowed; its kanji were attached later as kun'yomi (JMdict, KANJIDIC2).
\end{cousinweb}

\begin{grammarlens}[title={What you've built}]
One more word for everyday life at A2.
\end{grammarlens}

\subsection*{Your turn}
\begin{itemize}
\raggedright
  \item \emph{Say it:} \emph{omimai}
  \item \emph{Say it:} \emph{omimai}, once more
  \item \emph{Say it:} \emph{tomodachi no omimai ni ikimasen ka} — Shall we go and visit our sick friend?
  \item \emph{Recall:} say the Japanese for prevention, then the Japanese for to apply, then say \emph{omimai} again
\end{itemize}

\begin{tcolorbox}[breakable,colback=teal!4,colframe=teal!35!black,title={Before you move on}]
What does \ja{おみまい} mean? (\textquotedblleft{}A sick visit; calling on someone who is ill\textquotedblright{}.) Say it once more.
\end{tcolorbox}
\section[\texorpdfstring{\ja{かいふく} (kaifuku)}{かいふく (kaifuku)}]{\ja{かいふく} (kaifuku) — recovery}

\label{lesson:JA-C222-kaifuku}



\begin{quote}
Before the new one: say the Japanese for to apply, then the Japanese for a sick visit.
\end{quote}

\subsection*{You'll want to know: \ja{かいふく}}
\textbf{\ja{かいふく}} — \emph{kaifuku} — \textquotedblleft{}recovery\textquotedblright{}.

\begin{quote}
\textbf{\ja{かいふくは} \ja{はやいですか。}} — \emph{kaifuku wa hayai desu ka} — \textquotedblleft{}Is the recovery quick?\textquotedblright{}
\end{quote}

\begin{cousinweb}
\textbf{\ja{かいふく}} (\emph{kaifuku}) is Sino-Japanese (kango): two Chinese characters read with their on'yomi, readings borrowed from Chinese. They are \emph{kai}, \textquotedblleft{}to revolve, go round\textquotedblright{} (Middle Chinese \emph{hwoj}) and \emph{fuku}, \textquotedblleft{}to return\textquotedblright{} (\emph{bjuwk}) (KANJIDIC2; Middle Chinese in Baxter's transcription).
\end{cousinweb}

\begin{grammarlens}[title={What you've built}]
One more word for everyday life at A2.
\end{grammarlens}

\subsection*{Your turn}
\begin{itemize}
\raggedright
  \item \emph{Say it:} \emph{kaifuku}
  \item \emph{Say it:} \emph{kaifuku}, once more
  \item \emph{Say it:} \emph{kaifuku wa hayai desu ka} — Is the recovery quick?
  \item \emph{Recall:} say the Japanese for to apply, then the Japanese for a sick visit, then say \emph{kaifuku} again
\end{itemize}

\begin{tcolorbox}[breakable,colback=teal!4,colframe=teal!35!black,title={Before you move on}]
What does \ja{かいふく} mean? (\textquotedblleft{}Recovery\textquotedblright{}.) Say it once more.
\end{tcolorbox}
\section[\texorpdfstring{\ja{だいじょうぶ} (daijōbu)}{だいじょうぶ (daijōbu)}]{\ja{だいじょうぶ} (daijōbu) — all right; safe}

\label{lesson:JA-C222-daijoubu}



\begin{quote}
Before the new one: say the Japanese for a sick visit, then the Japanese for recovery.
\end{quote}

\subsection*{You'll want to know: \ja{だいじょうぶ}}
\textbf{\ja{だいじょうぶ}} — \emph{daijōbu} — \textquotedblleft{}all right; safe\textquotedblright{}.

\begin{quote}
\textbf{\ja{あしは} \ja{だいじょうぶですか。}} — \emph{ashi wa daijōbu desu ka} — \textquotedblleft{}Is your leg all right?\textquotedblright{}
\end{quote}

\begin{cousinweb}
\textbf{\ja{だいじょうぶ}} (\emph{daijōbu}) is Sino-Japanese (kango): three Chinese characters read with their on'yomi, readings borrowed from Chinese. They are \emph{dai}, \textquotedblleft{}big\textquotedblright{} (Middle Chinese \emph{dajH}), \emph{jō}, \textquotedblleft{}ten feet, a full height\textquotedblright{} (\emph{drjangX}) and \emph{bu}, \textquotedblleft{}man\textquotedblright{} (\emph{pju}) (KANJIDIC2; Middle Chinese in Baxter's transcription).
\end{cousinweb}

\begin{grammarlens}[title={What you've built}]
One more word for everyday life at A2.
\end{grammarlens}

\subsection*{Your turn}
\begin{itemize}
\raggedright
  \item \emph{Say it:} \emph{daijōbu}
  \item \emph{Say it:} \emph{daijōbu}, once more
  \item \emph{Say it:} \emph{ashi wa daijōbu desu ka} — Is your leg all right?
  \item \emph{Recall:} say the Japanese for a sick visit, then the Japanese for recovery, then say \emph{daijōbu} again
\end{itemize}

\begin{tcolorbox}[breakable,colback=teal!4,colframe=teal!35!black,title={Before you move on}]
What does \ja{だいじょうぶ} mean? (\textquotedblleft{}All right; safe\textquotedblright{}.) Say it once more.
\end{tcolorbox}
\section[\texorpdfstring{\ja{がまん} (gaman)}{がまん (gaman)}]{\ja{がまん} (gaman) — endurance; putting up with something}

\label{lesson:JA-C222-gaman}



\begin{quote}
Before the new one: say the Japanese for recovery, then the Japanese for all right.
\end{quote}

\subsection*{You'll want to know: \ja{がまん}}
\textbf{\ja{がまん}} — \emph{gaman} — \textquotedblleft{}endurance; putting up with something\textquotedblright{}.

\begin{quote}
\textbf{\ja{こどもは} \ja{がまんが} \ja{にがてです。}} — \emph{kodomo wa gaman ga nigate desu} — \textquotedblleft{}Children are bad at putting up with things.\textquotedblright{}
\end{quote}

\begin{cousinweb}
\textbf{\ja{がまん}} (\emph{gaman}) is Sino-Japanese (kango): two Chinese characters read with their on'yomi, readings borrowed from Chinese. They are \emph{ga}, \textquotedblleft{}self\textquotedblright{} (Middle Chinese \emph{ngaX}) and \emph{man}, \textquotedblleft{}ridicule, laziness\textquotedblright{} (\emph{maenH}) (KANJIDIC2; Middle Chinese in Baxter's transcription).
\end{cousinweb}

\begin{grammarlens}[title={What you've built}]
One more word for everyday life at A2.
\end{grammarlens}

\subsection*{Your turn}
\begin{itemize}
\raggedright
  \item \emph{Say it:} \emph{gaman}
  \item \emph{Say it:} \emph{gaman}, once more
  \item \emph{Say it:} \emph{kodomo wa gaman ga nigate desu} — Children are bad at putting up with things.
  \item \emph{Recall:} say the Japanese for recovery, then the Japanese for all right, then say \emph{gaman} again
\end{itemize}

\begin{tcolorbox}[breakable,colback=teal!4,colframe=teal!35!black,title={Before you move on}]
What does \ja{がまん} mean? (\textquotedblleft{}Endurance; putting up with something\textquotedblright{}.) Say it once more.
\end{tcolorbox}
